%=======================================================================
%  Paper V — Section 5 (What the method constrains)   FINAL
%
%  Requires the macros listed at the head of Section 2, plus:
%      \newcommand{\Dipole}{D}
%      \newcommand{\geomslope}{g}
%
%  PRINCIPAL FINDING OF THIS REVISION -------------------------------
%  Every statistic in ver10 of this section (A = 11.8, s = -1.14,
%  +/-13.1 %, +/-5.1 %) belongs to the ELEVEN planes at
%  x_split <= 0.2000 km, i.e. those below the LL-chondrite grain-density
%  ceiling. ver10 labelled this subset as the one the BLOCK-density
%  bound permits, which is the nine planes at x_split <= 0.1900 km and
%  gives different numbers. The eleven-plane subset spans a factor
%  1.907 in V_R, which is what the synthetic family's factor 1.90 was
%  matched to; the nine-plane subset spans 1.68 and all thirteen span
%  4.15. All three are now reported side by side in Table 5.1.
%
%  Superseded values:
%      "+/-13.1 % over the block-density range" -> 13.20 %, grain ceiling
%      "+/-0.435 %"                             -> 0.434 %
%      "factor of about ninety"                 -> decomposed, 3 subsets
%      "dipole displaced in observed direction" -> OPPOSITE direction
%      "no explanation for excess steepness"    -> Eq. (steepidentity)
%      "rho_rest = 1855 +/- 8"                  -> plane-placement spread;
%                                                  budget is +/- 27
%      "Delta M ... to +/-5.1 %"                -> +/-17.5 % with budget
%      dipole "three times less steady"         -> 21 sigma slope test
%
%  All values confirmed against paperV_final.json and synth_final.json.
%  Changes in this pass: the stability ratio is decomposed into identity,
%  subtraction and compensation (32.1 = 11.78 x 7.52 / 2.77), replacing the
%  single 'residual' of 2.7, which was the rho_R-vs-Delta M comparison under
%  another name; per-body synthetic table with standard errors; correlation
%  reported as Pearson with its two-tailed p and the Spearman value beside it.
%=======================================================================

\section{What the method constrains}
\label{sec:constrains}

Sections~\ref{sec:itokawa} and~\ref{sec:nulls} report one strong signal and
four weak ones. This section asks a different question: where the method
does respond strongly, what quantity has it determined? The answer is not
the density of the region, and establishing that is the principal result of
this paper.

\subsection{The observation, and the subset it depends on}
\label{sec:constrains:observation}

Across the thirteen-plane Itokawa sweep of Table~\ref{tab:itokawasweep} the
recovered region density varies by $\pm38.3\%$ about its mean while the
background density varies by $\pm0.434\%$. Both figures depend on how far
the sweep is allowed to run, and because the planes at the far end are ones
Section~\ref{sec:itokawa:ceiling} excludes on compositional grounds, the
dependence has to be made explicit before anything is inferred from it.

Table~\ref{tab:subsets} therefore reports every statistic over three nested
subsets: all thirteen planes; the eleven at
$\xsplit \le \SI{0.2000}{\kilo\metre}$, which lie below the LL-chondrite
grain-density ceiling and are the widest set that is compositionally
defensible; and the nine at $\xsplit \le \SI{0.1900}{\kilo\metre}$, which
lie below the softer empirical block-density bound adopted in
Section~\ref{sec:itokawa:ceiling}. We use the eleven-plane subset for the
scaling analysis of Sections~\ref{sec:constrains:slope}
to~\ref{sec:constrains:synthetic}, for one reason and not for its
convenience: it spans a factor $1.907$ in region volume, and the synthetic
family of Section~\ref{sec:constrains:synthetic} was built to span $1.90$.
A scaling exponent compared between sweeps of different span is not a
comparison. Where a conclusion changes between subsets we say so; the
central ones do not.

\begin{table}[htbp]
\centering
\small
\setlength{\tabcolsep}{5pt}
\caption{Sweep statistics over three nested subsets of the planes of
Table~\ref{tab:itokawasweep}. Coefficients of variation use the $n-1$
denominator. $\ampfac$ is evaluated per plane from
Equation~\eqref{eq:amplification} and averaged. The last three
ratios decompose $\CV(\densR)/\CV(\densrest)$ exactly: the identity of
Equation~\eqref{eq:excessidentity} fixes $\CV(\Dmass)/\CV(\densrest)$,
the subtraction that defines $\Ddens$ fixes $\CV(\Ddens)/\CV(\densR)$,
and only $\CV(\Ddens)/\CV(\Dmass)$, the compensation between contrast and
region volume, is a property of the objective function. The scaling exponent $s$ is
defined by Equation~\eqref{eq:slope} and fitted by ordinary least squares;
$\geomslope$ is the purely geometric slope of
Equation~\eqref{eq:geomslope}.}
\label{tab:subsets}
\begin{tabular}{lrrr}
\toprule
 & All planes & Below grain ceiling & Below block bound \\
$\xsplit$ range (km) & $0.1300$--$0.2400$ & $0.1300$--$0.2000$ & $0.1300$--$0.1900$ \\
$n$ & $13$ & $11$ & $9$ \\
span in $\VR$ & $4.15$ & $\mathbf{1.907}$ & $1.68$ \\
\midrule
$\densR$ mean (kg\,m$^{-3}$)      & $3169$   & $2936$   & $2856$ \\
\quad half-range                  & $38.26\%$ & $13.20\%$ & $10.07\%$ \\
\quad CV                          & $20.61\%$ & $8.60\%$  & $6.89\%$ \\
\midrule
$\densrest$ mean (kg\,m$^{-3}$)   & $1855.2$ & $1855.4$ & $1856.6$ \\
\quad half-range                  & $0.434\%$ & $0.434\%$ & $0.404\%$ \\
\quad CV                          & $0.251\%$ & $0.268\%$ & $0.252\%$ \\
\midrule
$\Dmass$ mean ($10^{9}$~kg)       & $2.805$  & $2.800$  & $2.779$ \\
\quad half-range                  & $5.12\%$  & $5.12\%$  & $4.80\%$ \\
\quad CV                          & $2.95\%$  & $3.16\%$  & $3.00\%$ \\
\midrule
$\dCOM$ mean (m)                  & $17.93$  & $17.51$  & $17.12$ \\
\quad half-range                  & $15.70\%$ & $12.82\%$ & $10.92\%$ \\
\quad CV                          & $9.31\%$  & $8.15\%$  & $7.42\%$ \\
\midrule
$\Ddens$ CV                       & $49.85\%$ & $23.79\%$ & $20.13\%$ \\
\midrule
$\bar{\ampfac}$                   & $11.77$  & $11.80$  & $11.90$ \\
$\CV(\densR)/\CV(\densrest)$      & $82.2$   & $32.1$   & $27.3$ \\
\quad identity, $\CV(\Dmass)/\CV(\densrest)$ & $11.76$ & $11.78$ & $11.88$ \\
\quad subtraction, $\CV(\Ddens)/\CV(\densR)$ & $2.42$ & $2.77$ & $2.92$ \\
\quad compensation, $\CV(\Ddens)/\CV(\Dmass)$ & $16.9$ & $\mathbf{7.52}$ & $6.71$ \\
\midrule
$s$                               & $-1.041\pm0.020$ & $\mathbf{-1.141\pm0.018}$ & $-1.159\pm0.025$ \\
$\geomslope$                      & $-0.177$ & $-0.244$ & $-0.262$ \\
$\mathrm{d}\log \Dipole/\mathrm{d}\log\VR$
                                  & $-0.217\pm0.020$ & $-0.385\pm0.018$ & $-0.420\pm0.025$ \\
\quad departure from dipole       & $10.9\sigma$ & $21.4\sigma$ & $16.8\sigma$ \\
\midrule
$\max \improve$                   & $26.4\%$ & $22.4\%$ & $21.0\%$ \\
\bottomrule
\end{tabular}
% Confirmed against paperV_final.json. The three values of geomslope were
% recomputed directly from the region centroids and x_cf = 0.36 m and agree.
\end{table}

Three things in that table carry the argument.

First, the background density is better determined than the region density
by a factor of $82$ over all thirteen planes and $32$ over the eleven. That
comparison is between $\densR$ and $\densrest$, which are measured in the
same units and are the two densities of the same model; they are not,
however, independent parameters, since the mass constraint of
Equation~\eqref{eq:massbalance} leaves exactly one degree of freedom, and
the pair $(\densR,\VR)$ is degenerate along $\Ddens\VR = \mathrm{const}$.
The inversion selects a point on that hyperbola only once the dividing
surface has been fixed by other means.

Second, only part of that factor is a result. Over the eleven planes it
decomposes exactly as $32.1 = 11.78 \times 7.52 / 2.77$.
Equation~\eqref{eq:excessamplify} fixes the first factor as algebra,
present whatever the objective function does. The third is also algebra:
$\Ddens$ is the difference between $\densR$ and a nearly constant
$\densrest$, so its coefficient of variation exceeds that of $\densR$ by the
ratio of their means. What remains, the factor of $7.5$ by which $\Dmass$ is
steadier than $\Ddens$, is the compensation between contrast and region
volume; since $\Dmass = \Ddens\VR$, it measures directly how far $\Ddens$
falls as $\VR$ grows, and it is the part we report. Dividing the net ratio
by $\bar{\ampfac}$ alone, as an earlier draft of this section did, gives
$\CV(\densR)/\CV(\Dmass) = 2.7$, which is the comparison of $\densR$ with
$\Dmass$ under another name: it leaves the subtraction factor in place and
understates the compensation by $2.8$.

Third, the ordering is not a property of the method. On 67P the region
density varies by $\pm3.1\%$ across its sweep while the excess mass varies
by $\pm34.4\%$---the reverse of Itokawa. On Arrokoth the density difference
$\Ddens$ passes through zero inside the sweep
(Section~\ref{sec:nulls:arrokoth}), so the sweep mean of $\Dmass$ lies
inside its own spread and no fractional range for it is meaningful; we
report instead that $\lvert\Dmass\rvert/\Mtot$ stays below $12\%$ at every
plane on both bodies. Figure~\ref{fig:rho_dm} shows the three sweeps on a
common scale.
These are half-ranges over the eight planes of the 67P sweep, from
$\xsplit = 0.70$ to $\SI{1.20}{\kilo\metre}$, which span a factor $1.504$ in
$\VR$ against the $1.907$ used for Itokawa.

\begin{figure}[p]
\centering
\includegraphics[height=0.86\textheight]{fig06_rho_vs_dM.pdf}
\caption{Recovered region density and excess mass across each split-plane
sweep, both as deviations from their own sweep mean and on a common
symmetric-log scale, against the region volume fraction. On Itokawa the
excess mass is nearly flat while the density slopes; on 67P and Arrokoth the
order is reversed. The shaded band marks $\pm20\%$. Arrokoth's excess mass
changes sign across its sweep, so its fractional range is not meaningful
and no scaling exponent exists. The Itokawa planes beyond
$\SI{0.2000}{\kilo\metre}$, which the grain-density ceiling of
Section~\ref{sec:itokawa:ceiling} excludes, are drawn open.}
\label{fig:rho_dm}
\end{figure}

\subsection{A diagnostic that does not break down}
\label{sec:constrains:slope}

Comparing fractional ranges fails when a quantity passes through zero,
because the mean then sits inside the spread. On the Arrokoth sweep
$\Ddens$ runs from $+25.0$ to $\SI{-61.5}{\kilo\gram\per\cubic\metre}$, so
its fractional range is meaningless; a ratio-based analysis of the synthetic
bodies produces apparent ratios of up to $1623$ that are artefacts of
dividing by nearly zero.

Since $\Dmass = \Ddens\,\VR$, the behaviour is captured instead by the
exponent
%
\begin{equation}
s \;=\; \frac{\mathrm{d}\log\lvert\Ddens\rvert}{\mathrm{d}\log \VR} ,
\label{eq:slope}
\end{equation}
%
fitted by ordinary least squares across the planes of a sweep. A slope of
$-1$ means $\Dmass$ is held fixed as the plane moves; a slope of $0$ means
the density difference is held fixed instead. Two conditions are required
for the exponent to be measurable: $\Ddens$ must keep its sign, and the
sweep must span an adequate range in $\VR$. We adopt a minimum span of
$1.35$, below which the logarithmic fit is unconstrained. The threshold
does no work in this paper: the Itokawa sweep spans $1.907$, the 67P sweep
$1.504$ and every synthetic sweep $1.90$ by construction, so no reported
exponent lies near it, and we report the standard error with every
exponent so that the reader can judge each fit directly.

Equation~\eqref{eq:excessidentity} gives the exponent a second reading, and
it is worth stating in its differential form rather than only at the single
point $s=-1$. Since $\Dmass = \Ddens\VR$ and
$\Dmass = (\densbulk - \densrest)\Vtot$ with $\densbulk$ and $\Vtot$ fixed,
%
\begin{equation}
s + 1 \;\equiv\; \frac{\mathrm{d}\log\Dmass}{\mathrm{d}\log\VR}
      \;\equiv\; \frac{\mathrm{d}\log\!\left(\densbulk-\densrest\right)}
                      {\mathrm{d}\log\VR} ,
\label{eq:steepidentity}
\end{equation}
%
of which the special case
%
\begin{equation}
s = -1 \iff \Ddens \VR = \mathrm{const} \iff \densrest = \mathrm{const}
\label{eq:slopeequiv}
\end{equation}
%
is the statement made in earlier drafts. Equation~\eqref{eq:steepidentity}
is an identity, not a result, and two consequences follow from it that
should be kept in view. The exponent is not independent evidence about
$\densrest$: it is a re-parametrization of the same information. What it
adds is direction rather than magnitude---it distinguishes a systematic
drift in the background density from scatter about a fixed value, which a
coefficient of variation cannot do---and it remains finite where fractional
ranges do not.

On the eleven Itokawa planes below the grain-density ceiling, $s = -1.141 \pm 0.018$. On 67P,
$s = -2.80 \pm 0.27$ over eight planes spanning a factor $1.504$ in $\VR$;
the uncertainty is fifteen times that on Itokawa, as expected on a body whose
dispersion improvement is $0.574\%$. On Arrokoth the density difference
changes sign and no exponent exists.

The Itokawa exponent is steeper than $-1$, and
Equation~\eqref{eq:steepidentity} says what that steepness is. It is the
systematic drift of the background density already documented in
Section~\ref{sec:itokawa:sweep}: $\densbulk - \densrest$ rises from $147.0$
to $\SI{163.1}{\kilo\gram\per\cubic\metre}$ as the region shrinks, giving
$\mathrm{d}\log(\densbulk-\densrest)/\mathrm{d}\log\VR = -0.16$ between the
endpoints, against the $-0.141$ that the fitted exponent implies. Earlier
drafts reported that the excess steepness was unexplained; it is not a
separate phenomenon from the shallow minimum in $\densrest$, and
Equation~\eqref{eq:steepidentity} makes the two the same statement.

\subsection{Excluding conservation of the mass dipole}
\label{sec:constrains:dipole}

A distinct hypothesis has to be disposed of before the excess mass can be
called the conserved quantity. Suppose the objective function held fixed not
$\Dmass$ but the mass dipole
$\Dipole = \Dmass\,(x_R - x_{\mathrm{cf}})$, which is $\Mtot\dCOM$ by
Equation~\eqref{eq:excessmass} and is the quantity the YORP measurement of
\citet{lowry2014} constrains. The two hypotheses make different and
quantitatively separated predictions, because the geometric factor is not
constant along a sweep.

Write
%
\begin{equation}
\geomslope \;\equiv\;
 \frac{\mathrm{d}\log\!\left(x_R - x_{\mathrm{cf}}\right)}{\mathrm{d}\log\VR} ,
\label{eq:geomslope}
\end{equation}
%
which involves no density and no objective function: it is fixed by the
shape model and the plane positions alone. Then
$\mathrm{d}\log \Dipole/\mathrm{d}\log\VR = (s+1) + \geomslope$, and dipole
conservation requires that combination to vanish, hence
$s = -1 - \geomslope$. On the eleven planes below the grain ceiling
$\geomslope = -0.244$, so dipole conservation predicts $s = -0.756$.

The measured value is $s = -1.141 \pm 0.018$. The prediction and the
measurement therefore lie on \emph{opposite} sides of $-1$: the observed
exponent is steeper, dipole conservation requires shallower. An earlier
draft of this section asserted the reverse, that dipole conservation would
displace the exponent in the observed direction, and that assertion was
wrong in sign.

Stated directly, the dipole itself scales as
%
\begin{equation}
\frac{\mathrm{d}\log \Dipole}{\mathrm{d}\log\VR}
 = -0.385 \pm 0.018 ,
\label{eq:dipoleslope}
\end{equation}
%
which is twenty-one standard errors from the zero that conservation
requires. The conclusion does not depend on the choice of subset: the same
calculation over all thirteen planes gives $-0.217 \pm 0.020$, and over the
nine planes below the block-density bound $-0.420 \pm 0.025$, or $10.9$ and
$16.8$ standard errors respectively (Table~\ref{tab:subsets}). A second and
cruder statement points the same way: across the eleven planes the
centre-of-mass offset has a coefficient of variation of $8.15\%$ against
$3.16\%$ for the excess mass, so the dipole is the less steady of the two
by a factor of $2.6$ rather than the more steady.
% Confirmed: geomslope recomputed directly from the region centroids and
% x_cf = 0.36 m reproduces -0.177, -0.244 and -0.262.


\subsection{A controlled experiment}
\label{sec:constrains:synthetic}

Three real bodies with three different behaviours do not settle why. A
fourth real body could not settle it either: we cannot choose in advance
whether it will carry a signal, and no body's true interior is known, so
agreement would establish only that something had happened twice.

We therefore constructed eight synthetic contact binaries as unions of
prolate lobes of revolution, meshed directly, with the strength of the
inversion signal varied by geometry alone. It should be stated at the outset
that no density anomaly is imposed: every synthetic body is of uniform
density by construction, and what varies between them is shape. The family
divides into three groups. A near-spherical control stands in for a body
with no interior structure. A symmetric family of equal lobes is symmetric
about its own neck, so the dispersion-minimizing contrast should be close to
unity and the response weak by construction. An asymmetric family of unequal
lobes has a geometric mass asymmetry, in the sense that a uniform body of
that shape carries its volume unequally about the dividing plane, and the
inversion has something to respond to. Within each family the neck is
deepened while the enclosed volume varies by no more than $7\%$.

Two design points matter. Planes are selected by bisection on the volume
fraction rather than by position, so that every body spans the same factor
$1.90$ in region volume as the eleven Itokawa planes below the grain-density ceiling, which span
$1.907$; selection by fixed distance produces spans of only $1.13$ to
$1.54$, which is not a comparison on equal terms. The full thirteen-plane
Itokawa sweep spans $4.15$ and is not comparable with the synthetic family
on this measure, which is the reason Table~\ref{tab:subsets} exists. And the
total mass is held at $\SI{2000}{\kilo\gram\per\cubic\metre}$ times the
enclosed volume with a rotation period of $\SI{12.13}{\hour}$ throughout, so
that only the shape differs.

It should be said clearly what this experiment does not test.
Potential-dispersion minimization does not fit observed data: it assumes the
true density is the one that flattens the surface geopotential. Imposing an
arbitrary density on a synthetic body and asking the method to recover it
would fail by construction and would test nothing. What is examined here is
the structure of the objective function---which quantity it holds fixed as
the dividing plane moves, and how that depends on the strength of its
response. Because no known density anomaly is injected, the experiment
characterises the regimes of the objective function; it does not by itself
establish that an exponent near $-1$ implies a physical mass anomaly in a
real body.

\begin{table}[htbp]
\centering
\small
\setlength{\tabcolsep}{5pt}
\caption{The synthetic family, one row per body. Improvements are the
maximum over each sweep, and the Itokawa figures given for comparison are
therefore also maxima---$22.4\%$ over the eleven planes below the grain
ceiling---rather than the $16.9\%$ reported at the detected neck in
Section~\ref{sec:itokawa:reference}. Every sweep uses seven planes spanning
a factor $1.90$ in $\VR$. Standard errors on $s$ are those of the
least-squares fit. $\bar{\ampfac}$ is the amplification factor of
Equation~\eqref{eq:excessamplify} averaged over the sweep; it is what
separates the $\densR$ and $\Dmass$ half-ranges, which are therefore not
directly comparable, and it is not meaningful on the three symmetric bodies,
where $\Dmass$ passes through zero. The last column is the compensation
ratio, the half-range of $\Ddens$ divided by that of $\Dmass$.}
\label{tab:synthetic}
\begin{tabular}{lrrrrrr}
\toprule
Body & $\improve_{\max}$ & $s$ & $\bar{\ampfac}$ & $\densR$ range
 & $\Dmass$ range & $\Ddens/\Dmass$ \\
\midrule
\texttt{sphere}        & $1.24\%$  & none & $44$ & $\pm3.10\%$ & sign change & --- \\
\texttt{sym\_waist}    & $3.68\%$  & none & ---  & $\pm5.61\%$ & sign change & --- \\
\texttt{sym\_neck}     & $4.35\%$  & none & ---  & $\pm6.22\%$ & sign change & --- \\
\texttt{sym\_deep}     & $4.47\%$  & none & ---  & $\pm5.60\%$ & sign change & --- \\
\addlinespace
\texttt{asym\_waist}   & $16.67\%$ & $-2.020\pm0.109$ & $15.0$ & $\pm9.28\%$  & $\pm31.3\%$ & $1.9$ \\
\texttt{asym\_neck}    & $17.63\%$ & $-2.366\pm0.165$ & $16.2$ & $\pm9.37\%$  & $\pm39.2\%$ & $1.7$ \\
\texttt{asym\_deep}    & $19.50\%$ & $-2.069\pm0.094$ & $14.9$ & $\pm8.81\%$  & $\pm29.1\%$ & $2.0$ \\
\texttt{asym\_extreme} & $43.71\%$ & $-1.217\pm0.023$ & $8.1$  & $\pm12.78\%$ & $\pm6.2\%$  & $6.1$ \\
\midrule
\itshape Itokawa, eleven planes & \itshape $22.4\%$ & \itshape $-1.141\pm0.018$
 & \itshape $11.8$ & \itshape $\pm13.20\%$ & \itshape $\pm5.12\%$
 & \itshape $7.1$ \\
\bottomrule
\end{tabular}
\end{table}

The pattern is monotonic. Where the response is weak the density difference
changes sign across the sweep and neither quantity is invariant, because
neither has anything to be invariant about. At moderate response the
exponent sits near $-2$, so the excess mass rises as the region shrinks and
the density is the steadier of the two. Only at the strongest response does
the exponent approach $-1$ and the excess mass become the steadier quantity.
Figure~\ref{fig:slope} plots the relation with the three real bodies
overlaid.

The Pearson correlation between signal strength and exponent is $0.95$
across the $n = 4$ synthetic bodies for which an exponent exists
($p = 0.05$ two-tailed, $0.03$ one-tailed); the four bodies without an
exponent cannot enter it. The Spearman rank correlation over the same four
is only $0.40$ ($p = 0.6$), because the three clustered bodies are not in
rank order: the Pearson value is carried by the single body at $43.7\%$. Four points, three of them clustered between $16.7\%$ and
$19.5\%$, cannot establish a functional form, and we return to this in
Section~\ref{sec:constrains:limits}.


\begin{figure}[htbp]
\centering
\includegraphics[width=\textwidth]{fig07_slope_vs_signal.pdf}
\caption{Scaling exponent $s=\mathrm{d}\log\lvert\Ddens\rvert/
\mathrm{d}\log\VR$ against signal strength for the synthetic family, with
the real bodies overlaid. A slope of $-1$ means the excess mass is held
fixed and $0$ that the density difference is. Bodies whose density
difference changes sign have no exponent and are placed on the lower row:
the symmetric synthetic family and Arrokoth. The exponent approaches $-1$
only as the signal strengthens, with a Pearson correlation of $0.95$ over the four
synthetic bodies that have one. Itokawa, at $-1.141\pm0.018$, is matched
only by the strongest synthetic body. 67P does not follow the pattern: its
exponent of $-2.80$ is steeper than the moderate group while its dispersion
improvement, $0.574\%$, places it among the weak bodies, which in the
synthetic family have no exponent at all.}
\label{fig:slope}
\end{figure}

The single synthetic body that reproduces Itokawa's behaviour is the one
with the strongest response, and it reproduces it on every measure:
exponent $-1.22$ against $-1.141$, density half-range $\pm12.78\%$ against
$\pm13.20\%$, excess-mass half-range $\pm6.22\%$ against $\pm5.12\%$. The
agreement is close on the eleven-plane subset and would not be on the
others---the thirteen-plane density half-range is $\pm38.3\%$---which is
the sharpest illustration of why the subset has to be declared before the
comparison is made.

One real body does not follow the pattern and we record it rather than
absorb it. 67P returns a measurable exponent of $-2.80$ while its
dispersion improvement, $0.574\%$, lies in the range where every synthetic
body has no exponent at all. Either the correspondence between improvement
and exponent is not one-to-one, or the 67P exponent is being measured on a
sweep whose $\Ddens$ is small enough throughout that the fit is dominated by
the systematics catalogued in Section~\ref{sec:nulls:67p}. We cannot
separate the two with the present data.

\subsection{The claim, restricted to the strong-signal regime}
\label{sec:constrains:claim}

Where the dispersion responds strongly to a contrast, the quantity the
inversion determines is the background density $\densrest$, and equivalently
through Equation~\eqref{eq:excessidentity} the mass anomaly $\Dmass$; the
pair $(\densR, \VR)$ is degenerate along $\Ddens \VR = \mathrm{const}$.
Where the response is weak, $\densrest$ drifts with the dividing plane and
nothing is determined. This accounts for why Itokawa behaves as it does and
why Arrokoth and the symmetric synthetic bodies do not, without appeal to
anything peculiar to any of them; 67P is the case it does not account for.

What the synthetic experiment establishes is a property of the objective
function in a regime, not a physical inference. It shows that when the
minimization responds strongly, the quantity it pins down is $\Dmass$ rather
than $\densR$. Whether a strong response on a real body \emph{implies} a
physical mass anomaly is a separate question that these data do not settle,
since no known anomaly was injected. On Itokawa the case for a physical
anomaly rests elsewhere: on the independent YORP determination of
\citet{lowry2014}, which makes no appeal to surface relaxation---and which,
by Section~\ref{sec:constrains:dipole}, constrains a quantity the objective
function demonstrably does not hold fixed, so the two lines of evidence are
genuinely independent rather than two readings of the same one.

The practical consequence is a reporting rule. Publish $\densrest$ and
$\Dmass$, which the method constrains, each with the amplification factor
$\ampfac$ and the plane at which it was evaluated; publish $\densR$ only
together with the dividing surface that produced it, since it is the ratio
of a constrained mass to a volume the investigator chose.

For Itokawa the numbers are these, and the distinction between the two kinds
of interval matters more than either value. Over the eleven planes below
the grain ceiling the background density has a mean of
$\SI{1855}{\kilo\gram\per\cubic\metre}$ with a plane-placement half-range of
$\SI{8}{\kilo\gram\per\cubic\metre}$, and the excess mass a mean of
$2.80\times10^{9}$~kg with a half-range of $5.1\%$. Those spreads describe
how little the answer moves when the plane moves; they are not
uncertainties, and quoting them as such---as an earlier draft of this
section did, in the form $1855\pm8$---understates the interval by more than
a factor of three. The uncertainty is the one assembled in
Equation~\eqref{eq:budget}, dominated by the curvature sensitivity interval:
$\densrest = 1858\pm\SI{27}{\kilo\gram\per\cubic\metre}$ and
$\Dmass = (2.75\pm0.48)\times10^{9}$~kg, or $7.7\pm1.3\%$ of the total, both
evaluated at the reference plane.

This also frames the head--neck ambiguity of
Section~\ref{sec:itokawa:ambiguity} correctly. A mass anomaly does not
specify the region that carries it. Two structures placing comparable excess
mass at comparable distance from the centre of figure will produce
comparable centre-of-mass offsets and comparable dispersion reductions, and
the method cannot separate them. That is not a defect of this implementation
but a property of what the data determine. It is also why the background
densities returned by the two region shapes differ by
$\SI{85}{\kilo\gram\per\cubic\metre}$: the invariance established here holds
under a change of plane position within a fixed region shape, and not under
a change of region shape.

\subsection{Limits of this evidence}
\label{sec:constrains:limits}

Five limits should be stated plainly.

The experiment varies shape, not density. Every synthetic body is uniform,
so what the exponent tracks is the strength of the inversion's response, and
the identification of that response with a physical mass anomaly is an
interpretation rather than a measured result. A version of this experiment
that built a body whose surface had relaxed onto an equipotential of a
planted interior would test the interpretation directly;
Section~\ref{sec:discussion:rate} explains why we do not report one.

The exponent is not independent of the background density.
Equation~\eqref{eq:steepidentity} makes $s+1$ identical to the logarithmic
drift of $\densbulk-\densrest$, so the scaling analysis and the stability
analysis of Section~\ref{sec:constrains:observation} are two readings of the
same data and should not be counted as two pieces of evidence. What the
exponent adds is the distinction between drift and scatter, and a
well-defined quantity where fractional ranges diverge. The dipole test of
Section~\ref{sec:constrains:dipole} is not subject to this objection,
because it compares the measured exponent against a value fixed by geometry
alone.

Four of the eight synthetic bodies produce a measurable exponent, and three
of those cluster between $16.7\%$ and $19.5\%$ improvement while the fourth
sits at $43.7\%$. The correlation of $0.95$ therefore rests on a single
high-leverage point, and with $n=4$ its nominal $p$ of $0.03$ should not be
read as strong. The \emph{direction} of the effect is supported; its
functional form is not determined, and a threshold anywhere between roughly
$20\%$ and $44\%$ fits the data equally well.

The numerical threshold does not transfer. The synthetic bodies require
about $44\%$ improvement to reach excess-mass invariance while Itokawa
reaches it at $22\%$ over the eleven planes below the grain ceiling, or $26\%$ if the
compositionally excluded planes are counted. They are smooth surfaces of
revolution; real bodies carry roughness, which Paper~IV found contributes an
error floor of order two percentage points. No threshold value should be
quoted as universal, and we quote none.

The experiment uses the same forward model in both directions. It therefore
examines the structure of the objective function, not the accuracy of the
forward model. That accuracy was established separately in Paper~IV against
an independent field, and is tested again here against an analytic ellipsoid
in Section~\ref{sec:discussion:ellipsoid}.

\subsection{Sensitivity to the assumed mass as a candidate diagnostic}
\label{sec:constrains:mass}

The mass-assumption tests of Sections~\ref{sec:itokawa:mass}
and~\ref{sec:nulls:arrokoth} were run for a different purpose, but placed
side by side they suggest a second, provisional way of separating a
strong-signal case from a weak one. The comparison has to be made carefully,
because the obvious reading of it is the wrong one.

\begin{table}[htbp]
\centering
\small
\caption{Response of the recovered contrast to the assumed total mass. The
final column is the logarithmic sensitivity
$\mathrm{d}\log C/\mathrm{d}\log\Mtot$, which is what makes the two rows
comparable: the mass ranges tested differ by a factor of two, so the raw
changes in the third column do not. On the raw change Itokawa appears the
\emph{more} sensitive of the two, which is the opposite of the intended
reading; per e-fold of assumed mass, Arrokoth is the more sensitive by a
factor of $1.6$. The property that separates the two bodies is not the size
of the change but whether the resulting interval reaches unity.}
\label{tab:massdiag}
\begin{tabular}{llllr}
\toprule
Body & Mass range & Contrast range & Behaviour
 & $\mathrm{d}\log C/\mathrm{d}\log\Mtot$ \\
\midrule
Itokawa  & factor $4$ & $1.404$--$1.569$ & excludes unity throughout
 & $0.080$ \\
Arrokoth & factor $2$ & $0.9159$--$0.9999$ & reaches unity
 & $0.127$ \\
\bottomrule
\end{tabular}
\end{table}

On the body with a strong response the inference is robust to the mass
assumption, the contrast staying clear of unity across a factor of four in
assumed mass; on the one without, a factor of two carries the answer to the
threshold of a sign reversal. The diagnostic is therefore the position of
the interval relative to unity and not its width. The test is cheaper than a
split-plane sweep---three or four runs at different assumed masses, with no
bisection and no plane search---though a dividing region is of course still
required.

We put this forward as a candidate diagnostic rather than an established
one, and with a specific reservation beyond the obvious ones. It rests on
two bodies that differ in several respects besides signal strength, among
them spin rate, size and whether the mass is measured at all, and those are
not separated here. The last of them is awkward: the diagnostic is of use
precisely where the mass is unknown, and on such a body there is no measured
value against which to calibrate the interval, while on a body where the
mass is known---Itokawa---the diagnostic is not needed. Whether it
generalises is untested, and the synthetic family of
Section~\ref{sec:constrains:synthetic}, where the mass is known exactly by
construction and the signal strength can be dialled, is the natural place to
test it.